% !TeX root = adelic-topology.tex
    The main reference of this section is \cite{Xie25ZDOsurfaces}.
    Fix an algebraically closed field \(\kkk\) which is the algebraic closure of a finitely generated field over \(\bbQ\).
    Let \(X\) be variety over \(\kkk\).

    In this section, we allow that the variety \(X\) is not irreducible,
    i.e., a variety over \(\kkk\) is a reduced separated scheme of finite type over \(\kkk\).

    Let \(L/K\) be a field extension.
    We denote by \(M_{L/K}\) (resp. \(M_{L,K})\) the set of places of \(L\) which restrict to \(K\) is trivial (resp. non-trivial).


\subsection{Adelic topology}

    \begin{notation}\label{notation:embedding_set_for_adelic_topology}
        Set
        \[ F(\kkk) \coloneqq \{\kk \subseteq \kkk \mid \kk \text{ is finitely generated over } \bbQ \text{ and } \kkk/\kk \text{ is algebraic}\}. \]
        Denote by \(I_\kk\) the set of all embeddings \(\sigma: \kk \to \bbC_\sigma\) over \(\kk\), where \(\bbC_\sigma = \bbC_p\) or \(\bbC\).
        Every such \(\sigma\) corresponds to a place \(v\in M_{\kk,\bbQ}\) by pulling back the standard absolute value on \(\bbC_p\) or \(\bbC\).
        For each \(\sigma\in I_\kk\), set \(E_\sigma = \{\bar{\sigma}: \kkk \to \bbC_\sigma \mid \bar{\sigma}|_\kk = \sigma\}\) be the set of embeddings of \(\kkk\) to \(\bbC_v\) extending \(\sigma\).

        We say that \(\kk\) is a \emph{defined field of \(X\)} if there exists a \(\kk\)-variety \(X_\kk\) such that \(X_\kk \times_{\kk} \Spec \kkk \cong X\).
        Such a variety \(X_\kk\) is called a \emph{model of \(X\) over \(\kk\)} or a \emph{\(\kk\)-model of \(X\)}.
        For every \(\sigma \in I_\kk\) and \(\bar{\sigma} \in E_\sigma\), we have an inclusion map \(\phi_{\bar{\sigma}}: X(\kkk) \cong X_{\kk}(\kkk) \to X_{\kk}(\bbC_{\sigma})\).
        Here the first isomorphism is given by the base change \(X_\kk \times_{\kk} \Spec \kkk \cong X\) and the second map is given by composition with \(\Spec \bar{\sigma}: \Spec \bbC_\sigma \to \Spec \kkk\).
        On \(X_{\kk}(\bbC_\sigma)\), we have the analytic topology induced from the norm on \(\bbC_\sigma\).
    \end{notation}

    \begin{definition}\label{def:basic_adelic_subset}
        Let \(\kk\) be a defined field of \(X\) with a model \(X_\kk\).
        Let \(\sigma,\sigma_i \in I_\kk\) and \(U \subset X_{\kk}(\bbC_\sigma), U_i \subseteq X_{\kk}(\bbC_{\sigma_i})\) be an open subsets in the analytic topology.
        We define
        \[ X_\kk(\sigma,U) \coloneqq \bigcup_{\tau\in E_\sigma} \phi_\tau^{-1}(U) \subseteq X(\kkk) \]
        and
        \[ X_\kk(\{\sigma_i,U_i\}_{i=1}^n) \coloneqq \bigcap_{i=1}^n X_\kk(\sigma_i,U_i) \subseteq X(\kkk). \]
        The subset of form \(X_\kk(\{\sigma_i,U_i\}_{i=1}^n)\) is called a \emph{basic adelic open subset} of \(X(\kkk)\).
    \end{definition}

    An equivalent description of \(X_\kk(\sigma,U)\) is the following.
    Given a field \(\kk \in F(\kkk)\) and a \(\kk\)-variety \(X_\kk\), fix an embedding \(\kk \injmap \kkk \injmap \bbC_\sigma\).
    There is a natural action of \(\Gal(\kkk/\kk)\) on \(X_\kk(\kkk)\), namely \((\tau,x) \mapsto \tau(x)\).
    For an analytic open subset \(U \subseteq X_{\kk}(\bbC_\sigma)\), we use the Galois action to ``spread'' \(\kkk\)-points in \(U\) to get a subset \(\bigcup_{\tau \in \Gal(\kkk/\kk)} \tau(U)\) of \(X_\kk(\kkk)\), which is exactly \(X_\kk(\sigma,U)\).
    Since we take the union over \(\Gal(\kkk/\kk)\), the subset \(X_\kk(\sigma,U)\) is independent of the choice of embedding \(\kkk \to \bbC_\sigma\).

    \begin{remark}\label{rmk:motivation_of_adelic_topology}
        When we consider dynamics on varieties, in particular, when we want to analyze the density of preperiodic points,
        the Zariski topology is too coarse to capture the density of points.
        For example, consider a non-automorphic dynamics \(f\) on \(\bbP^1\).
        It has infinitely many preperiodic points and hence they are always Zariski dense in \(\bbP^1\).
        However, if we consider the analytic topology on \(\bbP^1(\bbC)\), the preperiodic points sometimes are not dense.

        When we consider over number fields, we have no natural ``analytic topology'' on \(\bbP^1(\overline{\bbQ})\).
        An idea is to consider all possible places.
        That is, we can consider the topology generated by the union of all analytic topologies on \(\bbP^1(\bbC_v)\) for all places \(v\) of \(\overline{\bbQ}\).
        A good thing is that we have the Artin-Whaples Approximation Theorem, which says that the intersection of any finite collection of non-empty open subsets in \(\bbP^1(\bbC_v)\) for different places \(v\) is still non-empty.
        This tells us we still have some kind of irreducibility in this topology, which is very useful for the study of dynamics.
        However, the Artin-Whaples Approximation Theorem only works for different places.
        On the same place, it is clear that the intersection of two non-empty open subsets can be empty.

        In \cite{Xie25ZDOsurfaces}, the author defines the adelic topology.
        The key point is that we can take the union of all Galois conjugates of an open subset for a fixed place \(v\).
        Since after extending the base field, the same place \(v\) can split into several places, we can still have some kind of irreducibility on the same place, which is the \cref{lem:intersection_of_basic_adelic_subset_is_nonempty}.
    \end{remark}

    \begin{lemma}[{ref. \cite[Proposition 3.9]{Xie25ZDOsurfaces}, cf. \cite[Theorem 1.2]{ManZan14:AWapproxOnVarieties}}]\label{lem:intersection_of_basic_adelic_subset_is_nonempty}
        Suppose that \(X\) is irreducible.
        For any finite collection of basic adelic open subsets \(A_i\), the intersection \(\bigcap_{i=1}^n A_i\) is non-empty.
    \end{lemma}
    \begin{slogan}
        On a variety, the small open ball in one place will be dense with respect to other places.
    \end{slogan}
    \begin{proof}
        We follow the idea of \cite{ManZan14:AWapproxOnVarieties}.
        The strategy is to reduce to the case of curves by Bertini's Theorem and then apply the Riemann-Roch Theorem to the case of \(\bbP^1\).

        By passing to a finite extension of the base field, we can assume that for each basic adelic open subset \(A_i\),
        there is one place \(v_i \in M_\kkk\) and a non-empty analytic open subset \(U_i \subseteq X_{\kk}(\kkk)\) such that \(U_i \subset A_i\) and \(v_i\) are distinct for different \(i\).
        We only need to show that \(\bigcap_{i=1}^n U_i\) is non-empty.
        Since the singular locus of \(X\) is a proper Zariski closed subset, we can assume that \(X\) is smooth by replacing \(X\) with the smooth locus of \(X\).
        Furthermore, we can assume that \(X\) is affine by replacing \(X\) with an affine open subset of \(X\).
        Choose \(x_i \in U_i\) for each \(i\).
        By Bertini's Theorem, we can find a curve \(C\) in \(X\) passing through all \(x_i\).
        Replacing \(X\) by \(C\), we can assume that \(X\) is a curve.
        Again, we can assume that \(X\) is smooth and projective by replacing \(X\) with its normalization and then its smooth projective model.
        If necessary, we can re-choose \(x_i\) in \(U_i\).

        We do induction on \(n\).
        Suppose \(n=2\).
        By the Riemann-Roch Theorem, we can find a non-constant rational function \(f\) on \(X\) with only pole at \(x_1\).
        Then \(f\) induces a morphism \(f: X \to \bbP^1\) which is totally ramified at \(x_1\), i.e., \(f^{-1}(f(x_1)) = \{x_1\}\).
        By the Artin-Whaples Approximation Theorem, we can find \(y \in \bbP^1(\kkk)\) sufficiently closed to \(f(x_i)\) under \(v_i\) for \(i=1,2\) simultaneously.
        Take \(z \in f^{-1}(y) \cap U_2\), then \(z \in U_1 \cap U_2\) since \(f\) is totally ramified at \(x_1\) and \(y\) is sufficiently close to \(f(x_1)\) under \(v_1\).
        This finishes the case \(n=2\).\Yang{To be revised.since \cref{lem:image_sufficiently_closed_implies_preimage_sufficiently_closed}}

        For the general case, there exists a point \(z \in \bigcap_{i=1}^{n-1} U_i\) by induction hypothesis.
        By the Riemann-Roch Theorem, we can find a non-constant rational function \(f\) on \(X\) with only pole at \(z\).
        Then \(f\) induces a morphism \(f: X \to \bbP^1\) which is totally ramified at \(z\).
        By the Artin-Whaples Approximation Theorem, we can find \(y \in \bbP^1(\kkk)\) sufficiently close to \(f(z)\) under \(v_1,\ldots,v_{n-1}\) and sufficiently close to \(f(x_n)\) under \(v_n\) simultaneously.

        \Yang{To be added.}
    \end{proof}

    \begin{lemma}\label{lem:image_sufficiently_closed_implies_preimage_sufficiently_closed}
        Let \(f: X(\bbC_v) \to \bbP^1(\bbC_v)\) be a non-constant morphism.
        Fix any metric on \(\bbP^1(\bbC_v)\) and \(X(\bbC_v)\) compatible with the analytic topology.
        Let \(x \in X(\bbC_v)\) and \(y = f(x)\).
        Then for any \(\epsilon > 0\), there exists \(\delta > 0\) such that if \(y' \in \bbP^1(\bbC_v)\) with \(d(y,y') < \delta\), then there exists \(x' \in f^{-1}(y')\) with \(d(x,x') < \epsilon\).
    \end{lemma}

    \Yang{Flat morphism is open w.r.t. analytic topo.}

    Let \(\kk' \in F(\kkk)\) containing \(\kk\) and \(X_{\kk'} = X_\kk \times_{\kk} \Spec \kk'\).
    Then \([\kk':\kk] < +\infty\) and there are only finitely many \(\tau_1,\ldots,\tau_m \in I_{\kk'}\) extending \(\sigma\).
    Hence
    \[ X_\kk(\sigma,U) = \bigcup_{j=1}^m X_{\kk'}(\tau_j,U) \]
    for \(U \subseteq X_{\kk}(\bbC_\sigma) = X_{\kk'}(\bbC_\sigma)\).
    More generally, for any finite collection of \(\sigma_i \in I_\kk\) and \(U_i \subseteq X_{\kk}(\bbC_{\sigma_i})\), we have
    \[ X_\kk(\{\sigma_i,U_i\}_{i=1}^n) = \bigcap_{i=1}^n \left( \bigcup_{j=1}^{m_i} X_{\kk'}(\tau_{ij},U_i) \right) = \bigcup_{(j_1,\ldots,j_n)} X_{\kk'}(\{\tau_{ij_i},U_i\}_{i=1}^n). \]
    In conclusion, every basic adelic subset defined over a field \(\kk\) can be written as a finite union of basic adelic subsets defined over a finite extension \(\kk'\) of \(\kk\).

    Let \(f: Y \to X\) be a morphism of varieties over \(\kkk\) with a model \(f_\kk: Y_\kk \to X_\kk\) defined over \(\kk\).
    Given \(\sigma \in I_\kk\), we have an analytic map \(f_{\kk,\sigma}: Y_{\kk}(\bbC_\sigma) \to X_{\kk}(\bbC_\sigma)\) induced by \(f_\kk\).
    \Yang{}

    % \begin{lemma}\label{lem:basic_fact_of_basic_adelic_subset}
    %     We have the following basic facts of basic adelic subsets:
    %     \begin{enumerate}
    %         \item let \(\kk'\) be a finite extension of \(\kk\) and \(X_{\kk'}(\{\sigma_i',U_i'\}_{i=1}^m)\) be a basic adelic subset of \(X(\kkk)\) defined by \(\kk'\);
    %         \item let \(f: Y \to X\) be a morphism of varieties over \(\kkk\) and \(X_\kk(\{\sigma_i,U_i\}_{i=1}^n)\) be a basic adelic subset of \(X(\kkk)\) defined by \(\kk\),
    %             then \(f^{-1}(X_\kk(\{\sigma_i,U_i\}_{i=1}^n)) = Y_\kk(\{\sigma_i,f^{-1}(U_i)\}_{i=1}^n)\);
    %         \item let \(\tau \in \Gal(\kkk/\kk)\) and \(X_{\kk'}(\{\sigma_i',U_i'\}_{i=1}^m)\) be a basic adelic subset of \(X(\kkk)\) defined by \(\kk'\), then \(\tau(X_{\kk'}(\{\sigma_i',U_i'\}_{i=1}^m)) = X_{\kk'}(\{\tau\sigma_i',\tau(U_i')\}_{i=1}^m)\);
    %     \end{enumerate}
    %     \Yang{}
    % \end{lemma}


    \begin{definition}\label{def:general_adelic_subset}
        A \emph{general adelic subset} of \(X(\kkk)\) is defined to a subset of the form \(\pi(B)\) with \(\pi: Y \to X\) a flat morphism of varieties and \(B\) a basic adelic open subset of \(Y(\kkk)\).
    \end{definition}

    \begin{remark}\label{rmk:general_adelic_subset_and_define_field}
        To define a general adelic subset of \(X(\kkk)\), there are two fields involved: the field \(\kk\) over which the basic adelic subset is defined,
        and the field \(\bfl\) over which the morphism \(\pi:Y \to X\) is defined.

        If we fix a defined field \(\kk_0\) of \(X\), by \cite[Proposition 3.15 (ii) and (v)]{Xie25ZDOsurfaces}, we can always choose \(\bfl = \kk_0\) and \(\kk\) is a finite extension of \(\kk_0\).
    \end{remark}

    % \begin{remark}\label{rmk:irreducible_variety_and_finite_union_of_general_adelic_subset}
    %     Note that our definition of general adelic subset is slightly different from the one in \cite{Xie25ZDOsurfaces}, where the author does not require that the variety is irreducible.
    %     % Note that on \cite{Xie25ZDOsurfaces}, the author does not require that varieties are irreducible,
    %     if so, the finite union of general adelic subsets will be general adelic subsets in this sense.
    %     \Yang{To be revised.}
    % \end{remark}

    \begin{lemma}\label{prop:finite_union_and_intersection_of_general_adelic_subset}
        The finite union and intersection of general adelic subsets are still general adelic subsets.
    \end{lemma}
    \begin{proof}
        Let \(B_1,B_2\) be two general adelic subsets of \(X(\kkk)\) defined by \(\pi_i: Y_i \to X\) and basic adelic subsets \(A_i \subseteq Y_i(\kkk)\) for \(i=1,2\).
        Let \(Y = Y_1 \times_X Y_2\) and \(\pi: Y \to X\) be the natural morphism.
        Since we work over a field of characteristic zero, \(Y\) is reduced.
        Hence \(\pi: Y \to X\) is a flat morphism between varieties.
        Denote by \(C_i = p_i^{-1}(A_i)\) with \(p_i: Y \to Y_i\) the natural projection for \(i=1,2\).

        \Yang{To be added.}
    \end{proof}


    By \cref{prop:finite_union_and_intersection_of_general_adelic_subset}, the general adelic subsets form a basis of topology on \(X(\kkk)\).
    Hence we have the following definition.

    \begin{definition}\label{def:adelit_topology}
        The \emph{adelic topology} on \(X(\kkk)\) is defined to be the topology generated by all general adelic subsets of \(X(\kkk)\),
        i.e., an adelic open subset is an arbitrary union of general adelic subsets.
    \end{definition}

    \begin{proposition}\label{prop:basic_properties_of_adelic_topology}
        We have the following properties of adelic topology:
        \begin{enumerate}
            \item adelic topology is finer than Zariski topology;
            \item morphisms of varieties are continuous with respect to adelic topology;
            \item flat morphisms of varieties are open with respect to adelic topology.
        \end{enumerate}
    \end{proposition}
    \begin{proof}
        \Yang{To be added.}
    \end{proof}

    \begin{proposition}\label{prop:galois_action_is_continuous_wrt_adelic_and_profinite_topo}
        Let \(\kk\) be a defined field of \(X\).
        The action of \(\Gal(\kkk/\kk)\) on \(X(\kkk)\), namely \((\sigma,x) \mapsto \sigma(x)\), is continuous with respect to the adelic topology on \(X(\kkk)\) and the profinite topology on \(\Gal(\kkk/\kk)\).
    \end{proposition}
    \begin{proof}
        \Yang{To be added.}
    \end{proof}

    Recall that on \(\bbP^1_\bbQ\), the Artin-Whaples Approximation Theorem says that for any finite collection of places \(v_1,\ldots,v_n\) of \(\bbQ\) and any open subsets \(U_i \subseteq \bbP^1(\bbQ_{v_i})\),
    the intersection \(\bigcap_{i=1}^n (U_i\cap \bbP^1(\bbQ))\) is non-empty.
    The following lemma is a generalization and it is the motivation of the definition of adelic subsets.

    \begin{theorem}\label{thm:intersection_of_adelic_open_subset}
        Adelic topology preserves the irreducibility of varieties.
        Explicitly, on a variety, the intersection of any finite collection of non-empty adelic open subsets is non-empty.
    \end{theorem}
    \begin{proof}
        \Yang{To be added.}
    \end{proof}



% \subsection{Adelic topology}

\subsection{Examples and relation to dynamics}



    \Yang{Is \(\bbA^1(\bbQ)\) adelic closed in \(\bbA^1(\overline{\bbQ})\) with adelic topology? If so, why?}

    \Yang{Describe all adelic closed subset in \(\bbA^1(\overline{\bbQ})\).}
